\section{AHSME 1969}
        \begin{enumerate}
        	\item (\textit{AHSME 1969}) When $x$ is added to both the numerator and denominator of the fraction 
$\frac{a}{b},a \ne b,b \ne 0$, the value of the fraction is changed to $\frac{c}{d}$. 
Then $x$ equals:


 $\text{(A) } \frac{1}{c-d}\quad \text{(B) } \frac{ad-bc}{c-d}\quad \text{(C) } \frac{ad-bc}{c+d}\quad \text{(D) }\frac{bc-ad}{c-d} \quad \text{(E) } \frac{bc+ad}{c-d}$


 

	\item (\textit{AHSME 1969}) If an item is sold for $x$ dollars, there is a loss of $15\%$ based on the cost. If, however, the same item is sold for $y$ dollars, there is a profit of $15\%$ based on the cost. The ratio of $y:x$ is:


 $\text{(A) } 23:17\quad \text{(B) } 17y:23\quad \text{(C) } 23x:17\quad \\ \text{(D) dependent upon the cost} \quad \text{(E) none of these.}$


 



 

	\item (\textit{AHSME 1969}) If $N$, written in base $2$, is $11000$, the integer immediately preceding $N$, written in base $2$, is:


 $\text{(A) } 10001\quad \text{(B) } 10010\quad \text{(C) } 10011\quad \text{(D) } 10110\quad \text{(E) } 10111$


 

	\item (\textit{AHSME 1969}) Let a binary operation $\star$ on ordered pairs of integers be defined by $(a,b)\star (c,d)=(a-c,b+d)$. Then, if $(3,3)\star (0,0)$ and $(x,y)\star (3,2)$ represent identical pairs, $x$ equals:


 $\text{(A) } -3\quad \text{(B) } 0\quad \text{(C) } 2\quad \text{(D) } 3\quad \text{(E) } 6$


 

	\item (\textit{AHSME 1969}) If a number $N,N \ne 0$, diminished by four times its reciprocal, equals a given real constant $R$, then, for this given $R$, the sum of all such possible values of $N$ is


 $\text{(A) } \frac{1}{R}\quad \text{(B) } R\quad \text{(C) } 4\quad \text{(D) } \frac{1}{4}\quad \text{(E) } -R$


 

	\item (\textit{AHSME 1969}) The area of the ring between two concentric circles is $12\frac{1}{2}\pi$ square inches. The length of a chord of the larger circle tangent to the smaller circle, in inches, is:


 $\text{(A) } \frac{5}{\sqrt{2}}\quad \text{(B) } 5\quad \text{(C) } 5\sqrt{2}\quad \text{(D) } 10\quad \text{(E) } 10\sqrt{2}$


 

	\item (\textit{AHSME 1969}) If the points $(1,y_1)$ and $(-1,y_2)$ lie on the graph of $y=ax^2+bx+c$, and $y_1-y_2=-6$, then $b$ equals:


 $\text{(A) } -3\quad \text{(B) } 0\quad \text{(C) } 3\quad \text{(D) } \sqrt{ac}\quad \text{(E) } \frac{a+c}{2}$


 

	\item (\textit{AHSME 1969}) Triangle $ABC$ is inscribed in a circle. The measure of the non-overlapping minor arcs $AB$, $BC$ and $CA$ are, respectively, $x+75^{\circ} , 2x+25^{\circ},3x-22^{\circ}$. Then one interior angle of the triangle is:


 $\text{(A) } 57\frac{1}{2}^{\circ}\quad \text{(B) } 59^{\circ}\quad \text{(C) } 60^{\circ}\quad \text{(D) } 61^{\circ}\quad \text{(E) } 122^{\circ}$


 

	\item (\textit{AHSME 1969}) The arithmetic mean (ordinary average) of the fifty-two successive positive integers beginning at 2 is:


 $\text{(A) } 27\quad \text{(B) } 27\frac{1}{4}\quad \text{(C) } 27\frac{1}{2}\quad \text{(D) } 28\quad \text{(E) } 28\frac{1}{2}$


 

	\item (\textit{AHSME 1969}) The number of points equidistant from a circle and two parallel tangents to the circle is:


 $\text{(A) } 0\quad \text{(B) } 2\quad \text{(C) } 3\quad \text{(D) } 4\quad \text{(E) } \infty$


 

	\item (\textit{AHSME 1969}) Given points $P(-1,-2)$ and $Q(4,2)$ in the $xy$-plane; point $R(1,m)$ is taken so that $PR+RQ$ is a minimum. Then $m$ equals:


 $\text{(A) } -\frac{3}{5}\quad \text{(B) } -\frac{2}{5}\quad \text{(C) } -\frac{1}{5}\quad \text{(D) } \frac{1}{5}\quad \text{(E) either }-\frac{1}{5}\text{ or} \frac{1}{5}.$


 

	\item (\textit{AHSME 1969}) Let $F=\frac{6x^2+16x+3m}{6}$ be the square of an expression which is linear in $x$. Then $m$ has a particular value between:


 $\text{(A) } 3 \text{ and } 4\quad \text{(B) } 4 \text{ and } 5\quad \text{(C) } 5 \text{ and } 6\quad \text{(D) } -4 \text{ and } -3\quad \text{(E) } -6 \text{ and } -5$


 

	\item (\textit{AHSME 1969}) A circle with radius $r$ is contained within the region bounded by a circle with radius $R$. The area bounded by the larger circle is $\frac{a}{b}$ times the area of the region outside the smaller circle and inside the larger circle. Then $R:r$ equals:


 $\text{(A) }\sqrt{a}:\sqrt{b} \quad \text{(B) } \sqrt{a}:\sqrt{a-b}\quad \text{(C) } \sqrt{b}:\sqrt{a-b}\quad \text{(D) } a:\sqrt{a-b}\quad \text{(E) } b:\sqrt{a-b}$


 

	\item (\textit{AHSME 1969}) The complete set of $x$-values satisfying the inequality $\frac{x^2-4}{x^2-1}>0$ is the set of all $x$ such that:


 $\text{(A) } x>2 \text{ or } x<-2 \text{ or} -1<x<1\quad  \text{(B) } x>2 \text{ or } x<-2\quad \\ \text{(C) } x>1 \text{ or } x<-2\qquad\qquad\qquad\quad \text{(D) } x>1 \text{ or } x<-1\quad \\ \text{(E) } x \text{ is any real number except 1 or -1}$


 

	\item (\textit{AHSME 1969}) In a circle with center $O$ and radius $r$, chord $AB$ is drawn with length equal to $r$ (units). From $O$, a perpendicular to $AB$ meets $AB$ at $M$. From $M$ a perpendicular to $OA$ meets $OA$ at $D$. In terms of $r$ the area of triangle $MDA$, in appropriate square units, is:


 $\text{(A) } \frac{3r^2}{16}\quad \text{(B) } \frac{\pi r^2}{16}\quad \text{(C) } \frac{\pi r^2\sqrt{2}}{8}\quad \text{(D) } \frac{r^2\sqrt{3}}{32}\quad \text{(E) } \frac{r^2\sqrt{6}}{48}$


 

	\item (\textit{AHSME 1969}) When $(a-b)^n,n\ge2,ab\ne0$, is expanded by the binomial theorem, it is found that when $a=kb$, where $k$ is a positive integer, the sum of the second and third terms is zero. Then $n$ equals:


 $\text{(A) } \frac{1}{2}k(k-1)\quad \text{(B) } \frac{1}{2}k(k+1)\quad \text{(C) } 2k-1\quad \text{(D) } 2k\quad \text{(E) } 2k+1$


 

	\item (\textit{AHSME 1969}) The equation $2^{2x}-8\cdot 2^x+12=0$ is satisfied by:


 $\text{(A) } log(3)\quad \text{(B) } \frac{1}{2}log(6)\quad \text{(C) } 1+log(\frac{3}{2})\quad \text{(D) } 1+\frac{log(3)}{log(2)}\quad \text{(E) none of these}$


 

	\item (\textit{AHSME 1969}) The number of points common to the graphs of
$(x-y+2)(3x+y-4)=0 \text{    and     } (x+y-2)(2x-5y+7)=0$
is:


 $\text{(A) } 2\quad \text{(B) } 4\quad \text{(C) } 6\quad \text{(D) } 16\quad \text{(E) } \infty$


 

	\item (\textit{AHSME 1969}) The number of distinct ordered pairs $(x,y)$ where $x$ and $y$ have positive integral values satisfying the equation $x^4y^4-10x^2y^2+9=0$ is:


 $\text{(A) } 0\quad \text{(B) } 3\quad \text{(C) } 4\quad \text{(D) } 12\quad \text{(E) } \infty$


 

	\item (\textit{AHSME 1969}) Let $P$ equal the product of 3,659,893,456,789,325,678 and 342,973,489,379,256. The number of digits in $P$ is:


 $\text{(A) } 36\quad \text{(B) } 35\quad \text{(C) } 34\quad \text{(D) } 33\quad \text{(E) } 32$


 

	\item (\textit{AHSME 1969}) If the graph of $x^2+y^2=m$ is tangent to that of $x+y=\sqrt{2m}$, then:


 $\text{(A) m must equal } \frac{1}{2}\quad \text{(B) m must equal  } \frac{1}{\sqrt{2}}\quad\\ \text{(C) m must equal } \sqrt{2}\quad \text{(D) m must equal } 2\quad\\ \text{(E) m may be any non-negative real number}$


 

	\item (\textit{AHSME 1969}) Let $K$ be the measure of the area bounded by the $x$-axis, the line $x=8$, and the curve defined by


 
 \[f={(x,y)\quad |\quad y=x \text{  when  } 0 \le x \le 5, y=2x-5 \text{  when  } 5 \le x \le 8}.\]


 Then $K$ is:


 $\text{(A) } 21.5\quad \text{(B) } 36.4\quad \text{(C) } 36.5\quad \text{(D) } 44\quad \text{(E) less than 44 but arbitrarily close to it}$


 

	\item (\textit{AHSME 1969}) For any integer $n>1$, the number of prime numbers greater than $n!+1$ and less than $n!+n$ is:


 $\text{(A) } 0\quad\qquad \text{(B) } 1\quad\\ \text{(C) } \frac{n}{2} \text{ for n even, } \frac{n+1}{2} \text{ for n odd}\quad\\ \text{(D) } n-1\quad \text{(E) } n$


 

	\item (\textit{AHSME 1969}) When the natural numbers $P$ and $P'$, with $P>P'$, are divided by the natural number $D$, the remainders are $R$ and $R'$, respectively. When $PP'$ and $RR'$ are divided by $D$, the remainders are $r$ and $r'$, respectively. Then:


 $\text{(A) } r>r' \text{  always}\quad \text{(B) } r<r' \text{  always}\quad\\ \text{(C) } r>r' \text{  sometimes and } r<r' \text{  sometimes}\quad\\ \text{(D) } r>r' \text{  sometimes and } r=r' \text{  sometimes}\quad\\ \text{(E) } r=r' \text{  always}$


 

	\item (\textit{AHSME 1969}) If it is known that $\log_2(a)+\log_2(b) \ge 6$, then the least value that can be taken on by $a+b$ is:


 $\text{(A) } 2\sqrt{6}\quad \text{(B) } 6\quad \text{(C) } 8\sqrt{2}\quad \text{(D) } 16\quad \text{(E) none of these}$


 

	\item (\textit{AHSME 1969}) 
                    \begin{figure}[H]
                        \centering
                        \includegraphics[width=0.5\textwidth]{Problem/AHSME_I_1969_26_Problem_Figure_1.png}
                    \end{figure}
                    


 A parabolic arch has a height of $16$ inches and a span of $40$ inches. The height, in inches, of the arch at the point $5$ inches from the center $M$ is:


 $\text{(A) } 1\quad \text{(B) } 15\quad \text{(C) } 15\frac{1}{3}\quad \text{(D) } 15\frac{1}{2}\quad \text{(E) } 15\frac{3}{4}$


 

	\item (\textit{AHSME 1969}) A particle moves so that its speed for the second and subsequent miles varies inversely as the integral number of miles already traveled. For each subsequent mile the speed is constant. If the second mile is traversed in $2$ hours, then the time, in hours, needed to traverse the $n$th mile is:


 $\text{(A) } \frac{2}{n-1}\quad \text{(B) } \frac{n-1}{2}\quad \text{(C) } \frac{2}{n}\quad \text{(D) } 2n\quad \text{(E) } 2(n-1)$


 

	\item (\textit{AHSME 1969}) Let $n$ be the number of points $P$ interior to the region bounded by a circle with radius $1$, such that the sum of squares of the distances from $P$ to the endpoints of a given diameter is $3$. Then $n$ is:


 $\text{(A) } 0\quad \text{(B) } 1\quad \text{(C) } 2\quad \text{(D) } 4\quad \text{(E) } \infty$


 

	\item (\textit{AHSME 1969}) If $x=t^{1/(t-1)}$ and  $y=t^{t/(t-1)},t>0,t \ne 1$, a relation between $x$ and $y$ is: 


 $\text{(A) } y^x=x^{1/y}\quad \text{(B) } y^{1/x}=x^{y}\quad \text{(C) } y^x=x^y\quad \text{(D) } x^x=y^y\quad \text{(E) none of these}$


 

	\item (\textit{AHSME 1969}) Let $P$ be a point of hypotenuse $AB$ (or its extension) of isosceles right triangle $ABC$. Let $s=AP^2+PB^2$. Then:


 $\text{(A) } s<2CP^2 \text{ for a finite number of positions of P}\quad\\ \text{(B) } s<2CP^2 \text{ for an infinite number of positions of P}\quad\\ \text{(C) } s=2CP^2 \text{ only if P is the midpoint or an endpoint of AB}\quad\\ \text{(D) } s=2CP^2 \text{ always}\quad\\ \text{(E) } s>2CP^2 \text{ if P is a trisection point of AB}$


 

	\item (\textit{AHSME 1969}) Let $OABC$ be a unit square in the $xy$-plane with $O(0,0),A(1,0),B(1,1)$ and $C(0,1)$. Let $u=x^2-y^2$, and $v=2xy$ be a transformation of the $xy$-plane into the $uv$-plane. The transform (or image) of the square is:


 
                    \begin{figure}[H]
                        \centering
                        \includegraphics[width=0.5\textwidth]{Problem/AHSME_I_1969_31_Problem_Figure_1.png}
                    \end{figure}
                    


 
                    \begin{figure}[H]
                        \centering
                        \includegraphics[width=0.5\textwidth]{Problem/AHSME_I_1969_31_Problem_Figure_2.png}
                    \end{figure}
                    


 


                    \begin{figure}[H]
                        \centering
                        \includegraphics[width=0.5\textwidth]{Problem/AHSME_I_1969_31_Problem_Figure_3.png}
                    \end{figure}
                    


 


                    \begin{figure}[H]
                        \centering
                        \includegraphics[width=0.5\textwidth]{Problem/AHSME_I_1969_31_Problem_Figure_4.png}
                    \end{figure}
                    


 


                    \begin{figure}[H]
                        \centering
                        \includegraphics[width=0.5\textwidth]{Problem/AHSME_I_1969_31_Problem_Figure_5.png}
                    \end{figure}
                    


 

	\item (\textit{AHSME 1969}) Let a sequence $\{u_n\}$ be defined by $u_1=5$ and the relationship $u_{n+1}-u_n=3+4(n-1), n=1,2,3\cdots.$If $u_n$ is expressed as a polynomial in $n$, the algebraic sum of its coefficients is:


 $\text{(A) 3} \quad \text{(B) 4} \quad \text{(C) 5} \quad \text{(D) 6} \quad \text{(E) 11}$


 

	\item (\textit{AHSME 1969}) Let $S_n$ and $T_n$ be the respective sums of the first $n$ terms of two arithmetic series. If $S_n:T_n=(7n+1):(4n+27)$ for all $n$, the ratio of the eleventh term of the first series to the eleventh term of the second series is:


 $\text{(A) } 4:3\quad \text{(B) } 3:2\quad \text{(C) } 7:4\quad \text{(D) } 78:71\quad \text{(E) undetermined}$


 

	\item (\textit{AHSME 1969}) The remainder $R$ obtained by dividing $x^{100}$ by $x^2-3x+2$ is a polynomial of degree less than $2$. Then $R$ may be written as:


 

$\text{(A) }2^{100}-1 \quad \text{(B) } 2^{100}(x-1)-(x-2)\quad \text{(C) } 2^{200}(x-3)\quad\\ \text{(D) } x(2^{100}-1)+2(2^{99}-1)\quad \text{(E) } 2^{100}(x+1)-(x+2)$


 

	\item (\textit{AHSME 1969}) Let $L(m)$ be the $x$ coordinate of the left end point of the intersection of the graphs of $y=x^2-6$ and $y=m$, where $-6<m<6$. Let $r=[L(-m)-L(m)]/m$. Then, as $m$ is made arbitrarily close to zero, the value of $r$ is:


 $\text{(A) arbitrarily close to } 0\quad\\ \text{(B) arbitrarily close to }\frac{1}{\sqrt{6}}\quad \text{(C) arbitrarily close to }\frac{2}{\sqrt{6}} \quad \text{(D) arbitrarily large} \quad \text{(E) undetermined}$


 

\end{enumerate}